\chapter{What a Protocol Means}
\label{ch:meaning}
\label{ch:grammar}

\index{global type}
\index{grammar}
A protocol is a grammar. Its terminal symbols are messages, each with a sender
and a receiver; its productions are sequence, choice, \texttt{par},
\texttt{join} and \texttt{rec}; the language it generates is the set of legal
conversations. A recorded run either is a sentence of that language or it is not.
Because every protocol reads as if it were the top-level one, requirements stack:
a larger protocol includes smaller ones by name.

This chapter is what LeanFM checks before anything runs: that the protocol can
be implemented at all, what each role owes, and which exchanges are calls.

\section{Two rules every protocol must keep}

\index{realizability}
\index{choice!race}
\textbf{One role decides each choice.} The decider is the role that sends the
first message of every branch. When two roles each start a branch, both could act
at once and the conversation would leave the protocol:

\begin{lstlisting}[language={},basicstyle=\ttfamily\footnotesize,breaklines=true]
protocol Race (Client, Shop) {
  choice {
    Cancel from Client to Shop;
  } or {
    Ship from Shop to Client;
  }
}
\end{lstlisting}


\begin{lstlisting}[language={},basicstyle=\ttfamily\footnotesize,breaklines=true]
no single role decides a choice, so it is a race: Client starts Cancel; Shop starts Ship
\end{lstlisting}


The fix is to pick one decider: the client asks to cancel and the shop answers,
or the shop ships and the client receives. In the literature this is the
non-local choice, or realizability, condition.

\index{projection}
\textbf{Every role that must know which branch was taken is told.} A role whose
later behaviour differs between branches has to receive something that tells
them apart. Here the bank has taken part, it holds the client's money, and in one
branch nobody tells it what happened:

\begin{lstlisting}[language={},basicstyle=\ttfamily\footnotesize,breaklines=true]
protocol Blind (Client, Shop, Bank) {
  h: Hold from Client to Bank;
  choice {
    Capture from Shop to Bank;
    Ship from Shop to Client;
  } or {
    SoldOut from Shop to Client;
  }
}
\end{lstlisting}


\begin{lstlisting}[language={},basicstyle=\ttfamily\footnotesize,breaklines=true]
Bank cannot tell which branch Shop chose
\end{lstlisting}


A role whose part has not started yet owes nothing to a branch that never
contacts it. In \texttt{PlaceOrder} the bank is never told that the shop sold
out, and that is fine: its part starts with \texttt{Charge}, which only one branch
sends. A service that handles each request on its own, such as a mailer, goes in
a nested \texttt{sessions} block and never needs to know the conversation around
it.

A third rule is about \texttt{cause}: a message may only answer something its
sender has sent or received. A role decides from its own history, never from
something it could not have seen.

\section{Projection: what each role owes}

\index{local type}
Projecting the protocol onto a role keeps the messages it sends or receives and
turns the rest into its view of the choices: \texttt{select} where it decides,
\texttt{offer} where another role does and the arriving message says which.
\texttt{PlaceOrder}, projected:

\begin{lstlisting}[language={},basicstyle=\ttfamily\footnotesize,breaklines=true]
-- Client
send CreateOrder to Shop;
offer {
  receive Quote from Shop;
  send Pay to Shop;
  offer {
    receive Receipt from Shop;
  } or {
    receive PaymentDeclined from Shop;
  }
} or {
  receive SoldOut from Shop;
}

-- Shop
receive CreateOrder from Client;
select {
  send Quote to Client;
  receive Pay from Client;
  send Charge to Bank;
  offer {
    receive Charged from Bank;
    send Receipt to Client;
  } or {
    receive Declined from Bank;
    send PaymentDeclined to Client;
  }
} or {
  send SoldOut to Client;
}

-- Bank
offer {
  receive Charge from Shop;
  select {
    send Charged to Shop;
  } or {
    send Declined to Shop;
  }
} or {
}
\end{lstlisting}


\index{typestate}
A local type is an API. Each point in the conversation is a state, and each state
allows exactly the moves the protocol allows there. The client's typestate:

\begin{lstlisting}[language={},basicstyle=\ttfamily\footnotesize,breaklines=true]
start : S0
S0.sendCreateOrder(m : CreateOrder) : S1    -- to Shop
S1.receive() : Quote -> S2 (from Shop) | SoldOut -> End (from Shop)
S2.sendPay(m : Pay) : S3    -- to Shop
S3.receive() : Receipt -> End (from Shop) | PaymentDeclined -> End (from Shop)
\end{lstlisting}


A client in \texttt{S1} can only receive; a client in \texttt{End} has nothing
left to do. Code generated from this (Chapter~\ref{ch:generated}) makes the wrong
move a type error.

When the protocol keeps both rules, roles that follow their local types get the
classic guarantees: no message arrives that its receiver cannot handle, the run
follows the protocol, and no role waits forever
\citep{honda2008multiparty}. Each holds under the theory's assumption that every
message arrives exactly once. Real channels lose, duplicate and reorder, so
LeanFM lets a protocol declare each channel's physics beside it; the TFTP chapter
(\ref{ch:tftp-physics}) shows a protocol that is perfect by the classic rules and
has a famous bug once its physics is declared.

\section{Calls}

\index{HTTP}
\index{call}
\texttt{cause} also tells LeanFM which messages are calls and which are their
returns. A message is the return of the nearest unanswered request the other way
in its cause chain; otherwise it is a request. \texttt{PlaceOrder}'s calls:

\begin{lstlisting}[language={},basicstyle=\ttfamily\footnotesize,breaklines=true]
o: CreateOrder  Client → Shop  POST /orders
  returns Quote 201 | SoldOut 409
p: Pay  Client → Shop
  returns Receipt 200 | PaymentDeclined 402
    c: Charge  Shop → Bank
      returns Charged | Declined
\end{lstlisting}


A \texttt{route} makes a request an HTTP call, and over HTTP every request gets a
response on its own connection. LeanFM checks that, and suggests the line to add
when it is missing:

\begin{lstlisting}[language={},basicstyle=\ttfamily\footnotesize,breaklines=true]
protocol Lookup (Browser, Api) {
  q: GetUser from Browser to Api route "GET /users/{id}";
}
\end{lstlisting}


\begin{lstlisting}[language={},basicstyle=\ttfamily\footnotesize,breaklines=true]
1. [http] GetUser (q) from Browser to Api never gets a response
   try: Answer it: `GetUserDone from Api to Browser cause q status 200;` (or the status it means).
\end{lstlisting}


A protocol that passes all of this is ready to say more: what is assumed about
the decisions nobody else can see, and what is promised.
